\section{Захист ключової інформації в оперативній пам’яті }

\subsection{Чому ключ, який надійно зберігається у зашифрованому файлі, може стати доступним у відкритому вигляді під час роботи програми?}

\subsection{Які ризики для ключової інформації створюють memory dump, core dump, crash dump та swap?}

\subsection{Для чого після використання секретних даних може виконуватися очищення відповідних ділянок пам’яті?}

\subsection{Чому гарантоване видалення секретної інформації з оперативної пам’яті може бути складнішим у мовах із автоматичним керуванням пам’яттю?}

\subsection{Чим з погляду захисту ключа відрізняється виконання криптографічної операції програмою та виконання цієї операції всередині HSM або іншого захищеного апаратного середовища?}